\section{Background and Related Work}
\subsection{Background}
% \todo{Add a figure of prefill + decode, and tranformer + FFN}
\begin{figure}[!tp]
  \centering
  \includegraphics[width=\linewidth]{Figures/Background/prefill-decode.pdf}
  \vspace{-15pt}
  \caption{The (a) \textit{prefill} and (b) \textit{decode} stage of LLMs. Colored squares are weights or cached data. Gray denotes activations.}
  \vspace{-15pt}
  \label{fig:prefill-decode}
\end{figure}

Transformer-based LLMs~\citep{touvron2023llama2, zhang2205opt} achieve state-of-the-art (SOTA) performance across all kinds of Natural Language Processing (NLP) tasks. 
The transformer model architecture consists of many cascaded transformer blocks and each transformer block generally includes two types of networks: the Multi-Head Attention (MHA) and the Feed Forward Network (FFN). 

Given $N$ input tokens embedded in $D$ dimensional space $X \in \mathbb{R}^{N \times D}$, the MHA projects the token embedding as $h$ heads' query $Q$, key $K$ and value $V$ matrices in $\mathbb{R}^{h \times N \times D}$ space and performs attention operation for each head as shown in equation~\ref{eq:mha}.

\begin{equation}
% \begin{aligned}
    Q = XW_{Q}, K = XW_{K}, V = XW_{V};  \\
    O = \text{softmax}(QK^{T})V
% \end{aligned}
\label{eq:mha}
\end{equation}
where $W_{Q}$, $W_{K}$, $W_{V}$ represent the weight matrix in MHA.

The FFN further transforms each token embedding. Given the MHA output matrix $O \in \mathbb{R}^{h \times N \times D}$, 
% , embeddings for each token from multiple heads are aggregated through linear transformation and get $O' \in \mathbb{R}^{N \times D}$. 
FFN passes it through two fully connected layers with a non-linear activation function $g$ and generates the token embedding for the next transformer block.

\begin{equation}
% \begin{aligned}
    X = g(O W_1)W_2
% \end{aligned}
\label{eq:ffn}
\end{equation}
where, $W_{1}$, $W_{2}$ represent the two weight matrix in FFN.

% \fty{we shall define the phrase. For example, the input from user is called context or request. The generated text from the model is called response.}
As shown in Fig.~\ref{fig:prefill-decode}, the workflow of LLMs can be divided into two main stages: the prefill stage and the decode stage.
% The text generation of LLMs is divided into two main stages: the prefill stage and the decode stage.
In the prefill stage, LLM takes a prompt from the user which is a sequence of tokens as the input (e.g. the "\textit{Who won ?}" in Figure.\ref{fig:prefill-decode} (a)).
Then, LLM will understand the context of the prompt and generates the first response token (e.g. the "\textit{Alex}" in Figure.\ref{fig:prefill-decode} (a)).
All the $N$ input tokens are processed simultaneously with high throughput.
% different layers of LLM compute the Key and Value of the current Prompt and cache them for subsequent computations. Since the input X used by different layers are the Embedding results corresponding to the current Prompt, the computation process in the prefill stage can be parallelised between different layers.
% The basic computation process is as follows:
In the decode stage, LLM treats the newly generated token as length $N=1$ input and generates the next token (e.g. the "\textit{won}" in Figure.\ref{fig:prefill-decode} (b)). Since LLM only processes one new token at a time in the decode stage, the matrix-matrix multiplications in equation~\ref{eq:mha} and ~\ref{eq:ffn} become matrix-vector multiplications.
% \fty{do we need a fig here to show attention and FFN, as well as prefill and decode?}. 
The decode stage is called iteratively to generate the response token by token.
% In this computation process, it is necessary to use the Key and Value obtained from the prefill stage of computation, so as to complete the computation of the Encoder within a single layer.
% Where K and V correspond to the key and value in the KV Cache, X is the input of the current layer, Q is the Query matrix obtained by computing X through the $W_Q$ matrix, and $\frac{Softmax(Q^{T}K)V}{\sqrt{d}}$ is also known as the Attention operation.
% FFN stands for Feed Forward Neural, which is used to transform the result of Attention's computation into features once more.
% \xhdr{Heavy Computation and memory access}

\vspace{-5pt}
\subsection{Related Work}
% \subsubsection{Efficient Transformer}
\textbf{Efficient Transformer.}
To tackle the extreme overhead of LLMs, various compression techniques are commonly used.
Quantization~\citep{dettmers2022llm_int8, frantar2022gptq, yao2022zeroquant, xiao2023smoothquant,lin2023awq} approaches use low-bit integers to substitute the 16-bit floating point parameters and activations for inference. 
% Unlike small-sized models, LLMs often generate outlier activations with large magnitude, leading to large quantization error and accuracy loss for trivial quantization methods~\citep{dettmers2022llm_int8}. 
% State-of-the-art solutions try to solve the outlier problem through weight-activation scale transformation~\citep{xiao2023smoothquant} and salient channel scaling~\citep{lin2023awq}, pushing the limit of weight bit width to 4 bits on GPUs.
% However, given the lack of mix-precision support of GPU, no previous work utilizes mix-precision quantization to reduce the bit width further to the best of our knowledge.
Sparse attention~\citep{child2019spTrans, zaheer2020bigbird, beltagy2020longformer, wang2021spatten, chen2023dynamic_nm} and weight pruning~\citep{frantar2023sparsegpt, kwon2022fast} approaches skip part of the attention matrix or weight matrix computing according to the defined sparse pattern. Various sparse patterns are proposed for different tasks and transformers, including local diagonal pattern~\citep{beltagy2020longformer}, block sparse~\citep{child2019spTrans, zaheer2020bigbird, kitaev2020reformer, kwon2022fast}, N:M sparse pattern~\citep{chen2023dynamic_nm, frantar2023sparsegpt}, row-column skipping pattern~\citep{wang2021spatten}, unstructured pattern~\citep{frantar2023sparsegpt} and so on.
% Previous work pushes the limit of attention and weight sparsity to 67\% (2 thousand tokens by~\citep{zaheer2020bigbird}) and 60\% (OPT-175B model by \citep{frantar2023sparsegpt}) or even higher. However, most of them fail to demonstrate corresponding speedup on GPU due to the lack of sparse computation support of the hardware~\citep{tay2020LRA}.

% \subsubsection{LLM-related Accelerators}
\textbf{LLM-related Accelerators.}
Previous work~\cite{li2020ftrans,ham2020a3,isca2021elsa,micro2021sanger,wang2021spatten,micro2022DFX,micro2022bufferfly,hpca23flat,isca2023fact,wang2023cosa} propose customized architecture design for transformer models. 
Some work~\cite{CTA2023, micro2021sanger,wang2021spatten,micro2022bufferfly,li2020ftrans,isca2021elsa} lay more emphasis on accelerating sparse attention.
They design specialized architectures to fully utilize the pre-defined static attention pattern~\cite{micro2022bufferfly,li2020ftrans} or dynamically generated attention pattern~\cite{isca2023fact,CTA2023,wang2021spatten,micro2021sanger}.
Recently, FACT~\cite{isca2023fact} points out the importance of compressing linear layers with mixed-precision quantization to help reduce latency.
However, these methods cannot accelerate the decode stage of LLMs since they mainly focus on the \textit{prefill} stage for discriminative models, like medium-sized BERT~\cite{devlin2019bert} model.
% SpAtten~\cite{wang2021spatten} reduces attention computation in both \textit{prefill} and \textit{decode} stage by directly pruning attention heads. 
DFX~\cite{micro2022DFX} emphasizes the acceleration of the decode stage of LLMs. However, it lacks hardware support for model compression of LLMs, making it hard to further expand model size or maximum token size.
